\subsection{Auxiliary source frames}
\label{sec:scratch-source-frames}

The transfer construction assigns a rational matrix to every source,
gate and sink, common to all incidences of a gate, and requires
\[
 M_{\mathrm{sink}(\rho(w))}-M_{\mathrm{source}(w)}=I_m
\]
for every input role $w$. The retained networks give every auxiliary
source the zero matrix and every auxiliary sink the identity. The transfer
does not need this choice.

\begin{lemma}[Auxiliary source frames]
\label{lem:scratch-source-frames}
Let $w$ be an auxiliary role, so that $\rho(w)=w$, and let $M$ be a fixed
rational $m\times m$ matrix. Assign $M$ to the source of $w$ and $I_m+M$
to its sink, and leave every other matrix unchanged. Then the scalar
network, every common gate matrix and every endpoint identity are
preserved, and the address map implemented on every role is unchanged.
Only the first and last edges of $w$ change.
\end{lemma}

\begin{proof}
The endpoint identity for $w$ reads $(I_m+M)-M=I_m$; no other endpoint
identity involves this source or sink. The physical input stream $f_w$ has
the logical interpretation $g_w=\Phi_{-M}f_w$. This is a definition, not a
tape operation, and $g_w$ is arbitrary because $f_w$ is. The scalar network
restores arbitrary auxiliary values, so the logical value at the sink is
again $g_w$. The common-frame identity gives the physical output
\[
 \Phi_{I_m+M}\Phi_{-M}f_w=\Phi_{I_m}f_w,
\]
the same shear as on every other role. Gates on $w$ keep their common
matrices, so every interior edge of $w$, and every edge of every other
role, is unchanged.
\end{proof}

The lemma moves rank between the two end edges of an auxiliary role.
Before batching only the total mattered. With contiguous blocks, a large
edge costs much less per rank unit than a small one, so concentrating a
role's rank in one edge lowers the cost.

\paragraph{Stage-two auxiliary roles.}
Use the notation of Section~\ref{sec:batched-bit-rank-accounting}.
At tensor stage two put $A=F$, $P=\langle t_{a_1}\rangle$, $B=P^\perp$ and
$Q=\langle t_{a_3}\rangle$. Every auxiliary role of a stage-two invocation,
including its central roles, has first gate frame
$D_0=(B\otimes F)\otimes Q$ of dimension $h^2-h$, and last gate frame
$D_1=(A\otimes F)\otimes Q$ of dimension $h^2$. Assign every such role the
source matrix $P_{D_0}$, and therefore the sink matrix $I+P_{D_0}$. Then:
\begin{enumerate}
\item the entrance edge has matrix $P_{D_0}-P_{D_0}=0$ and makes no call;
\item every interior edge is unchanged;
\item since $D_0\subset D_1$ are nondegenerate, $D_0\perp D_1^\perp$, and
the exit edge has matrix
\[
 I+P_{D_0}-P_{D_1}=P_{D_1^\perp}+P_{D_0}=P_{D_0\perp D_1^\perp}.
\]
This is an idempotent of rank $m-h$. Its nullspace is
$D_1\cap D_0^\perp=P\otimes F\otimes Q$, of dimension $h$, because
$A=B\perp P$.
\end{enumerate}
Formerly the entrance had rank $h^2-h$ and the exit rank $m-h^2$. Now the
entrance has rank zero and the exit rank $m-h$. The total rank of each role,
the rank sum $s$ and the deficit $Wm-s$ are therefore unchanged. The
stage-two sink class $A_2$ of Section~\ref{sec:controlled-projector-basis}
no longer occurs. Its replacement is
\[
 A_4=I-\Pi_{P\otimes F\otimes Q},
\]
where $\Pi_U$ denotes the orthogonal projector onto $U$ for the label form.

\begin{lemma}[A common basis for the new exit class]
\label{lem:source-frame-basis}
In the parameter family $S=T(I_F\otimes K)$ of
Lemma~\ref{lem:controlled-projector-basis}, there are rational
$K,G_1,\ldots,G_H$ satisfying its first two assertions and, in addition:
for every $A_4$, the submatrix of $SA_4S^{-1}$ in its first $h$ rows and
last $h$ columns is nonsingular.
\end{lemma}

\begin{proof}
Order coordinates as in Section~\ref{sec:controlled-projector-basis}: the
original third factor $F$ is slow and $V$, the first two factors, is fast.
Then the nullspace of $A_4$ is $\langle t\rangle\otimes X$ with
$t=t_{a_3}$ and $X=\langle t_{a_1}\rangle\otimes F\subset V$, of dimension $h$.
Write $\Pi=UV'$ with $V'U=I_h$. Its first-$h$-row by last-$h$-column block
after conjugation is the product of the first $h$ rows of $SU$ and the last
$h$ columns of $V'S^{-1}$. The corresponding block of $SA_4S^{-1}$ is its
negative, because these row and column index sets are disjoint. It is
therefore enough that both square factors can be nonsingular. Each is then a
polynomial condition that is not identically zero.

A nullspace vector $t\otimes x$, with $x\in X$, has coordinate
$(G_it)_\alpha(Kx)_i$ at $(\alpha,i)$ after applying $S$. The first $h$
coordinates have $\alpha=1$ and $1\le i\le h$, so the first factor is the
map $x\mapsto((G_it)_1(Kx)_i)_{i\le h}$. Give every $G_i$ the first row
$(1,\ldots,1)$, so that $(G_it)_1=|a_3|=5$, and choose $K$ sending a basis
of $X$ to $e_1,\ldots,e_h$. The map is then nonsingular.

The rows of $V'$ span the form dual $\Omega(\langle t\rangle\otimes X)
=\langle\Omega_Ft\rangle\otimes\Omega_VX$, where $\Omega$ is the tensor label
form. This subspace has the same shape. Since
$S^{-\mathsf T}=T^{-\mathsf T}(I_F\otimes K^{-\mathsf T})$ and
$T^{-\mathsf T}$ acts by $G_i^{-\mathsf T}$, the same argument at the last
$h$ coordinates, with $\alpha=h$ and $H-h<i\le H$, shows that the second
factor can be nonsingular.

As in Lemma~\ref{lem:controlled-projector-basis}, clear denominators and
multiply these nonzero polynomials over the finitely many $A_4$. Include
the polynomials for its first two assertions and the invertibility of
$K$ and every $G_i$. A nonzero rational polynomial does not vanish at every
rational point. Its third assertion concerned $A_2$ and is no longer needed.
\end{proof}

By the large-projector factorization of Section~\ref{sec:projector-batching},
every exit edge $A_4$ is compiled as $h$ singleton interchanges and one
contiguous interchange of its middle $m-2h$ fields. The first and third
classes of Section~\ref{sec:batched-bit-rank-accounting} keep their
singleton corners and contiguous middle blocks.

\subsection{The interchange rank moment with source frames}
\label{sec:source-frame-bit-moment}

Keep the retained counts
\[
 h=28,\quad v=\binom{28}{5}=98280,\quad m=21952,\quad R=11840940,
\]
\[
 W=2v^2(v+R),\qquad s=Wm-v^3+2L,
\]
and put $B=v^2R=114371146876896000$, the number of stage-two auxiliary
roles. The large classes are now
\[
\begin{array}{c|r|r|r|r}
 \text{class}&\text{copies}&\text{rank}&\text{singleton pivots}&
   \text{batched width}\\\hline
 \text{stage 1--3 join}&B&21896&56&21840\\
 \text{stage 3 data entrance}&2v^3&21141&811&20330\\
 \text{stage 2 auxiliary exit}&B&21924&28&21896
\end{array}
\]
Every other rank unit is a singleton interchange. Relative to the
controlled count $S_1$ of Section~\ref{sec:controlled-projector-basis},
each stage-two auxiliary role loses its entrance of $h^2-h$ singletons and
gains $h$ exit corner pivots. Its former sink corner and middle blocks are
replaced by the new exit. Hence
\[
 S'=S_1-B(h^2-2h)=22293445170300595200.
\]
The normalized child widths are
\[
 r=\left(\frac1{21952},\ \frac{195}{196},\ \frac{10165}{10976},\
       \frac{391}{392}\right),
\]
and their rank-mass weights are
\[
 w=\left(\frac{2289741}{520019360},\ \frac{12827685}{26000968},\
       \frac{20865}{2736944},\ \frac{25721153}{52001936}\right).
\]
The weights sum to $1-\eta$ with $\eta=39/520019360$, so batching again
preserves the rank deficit. The positive logarithm series used before give
\[
 \log m<\frac{9997}{1000},\quad
 \log\frac{196}{195}<\frac{512}{10^5},\quad
 \log\frac{10976}{10165}<\frac{768}{10^4},\quad
 \log\frac{392}{391}<\frac{2555}{10^6}.
\]
Put
\[
 a_{\rm b}=\frac{154}{10^8},\qquad \tau=1-a_{\rm b}.
\]
With $e^x\le(1-x)^{-1}$ for $0\le x<1$, exact rational arithmetic gives
\[
 \Psi(1-a_{\rm b})<\sum_{i=0}^3\frac{w_i}{1-a_{\rm b}\ell_i}<1,
\]
with gap
\[
 \frac{9704597579224556535778749633448639054699}
      {20663392005574270637817582110595278287793980851574}
 >4.69\times10^{-10}.
\]
The arbitrary-width recurrence of Section~\ref{sec:mixed-width-rows}
therefore gives the interchange saving $a_{\rm b}=154/10^8$ for the
unchanged PR~\#7 scalar producer.
